% Notas de aula — Capítulo 23 (módulo extra): Mudança Climática
% Curso: Meteorologia Prática (baseado em R. Stull, Practical Meteorology, CC BY-NC-SA 4.0)
% Caixas verdes = conteúdo literal do quadro; linhas verdes = encenação.
\documentclass[11pt]{article}
\usepackage{../comum/preambulo}
\graphicspath{{../figuras/cap23/}}

\title{Capítulo 23 (módulo extra) --- Mudança Climática\\
{\large Notas de aula (roteiro de quadro-negro)}}
\author{Curso de Meteorologia Prática}
\date{}

\begin{document}
\maketitle
\thispagestyle{fancy}

\noindent\textbf{Plano sugerido:} 2 aulas de 2\,h.
\emph{Aula 1:} \S\ref{sec:forcante}--\S\ref{sec:inercia} (forçante do
CO$_2$, sensibilidade, a inércia oceânica).
\emph{Aula 2:} \S\ref{sec:tcre}--\S\ref{sec:sintese} (orçamento de
carbono, detecção e atribuição, impactos e o fecho do curso).

\section{A física que o curso inteiro preparou}\label{sec:forcante}

\begin{noquadro}[abertura]
\textbf{Cap.~23 --- Mudança Climática}\\[2pt]
não é um tema novo: é o curso inteiro aplicado a uma perturbação\\
estufa (Cap.~2) $+$ feedbacks (Cap.~21) $+$ oceano lento (Cap.~3)
$+$ estatística vs.\ caos (Cap.~20)
\end{noquadro}

\subsection{A forçante logarítmica}

\quadro{Perguntar: se a banda do CO$_2$ já é opaca
\javimos{2}{$\tau\gg1$}, por que mais CO$_2$ importa? (Deixar a
turma se enrolar; a resposta são as ASAS da banda.)}

\begin{formadif}[o logaritmo (T1; Caso~1; N1)]
Centro saturado; asas com $k \sim e^{-|\nu-\nu_0|/\Delta}$: a
fronteira $\tau = 1$ avança $\propto \ln C$:
\[ \Delta F = 5{,}35\,\ln\frac{C}{C_0}\ \mathrm{W/m^2},\qquad
\Delta F_{2\times} = 3{,}7\ \mathrm{W/m^2}. \]
Cada duplicação compra as mesmas asas. Hoje ($\sim$420\,ppm):
$\Delta F \simeq 2{,}2$\,W/m² — medível por satélite no espectro de
saída \javimos{8}{sondadores}.
\end{formadif}

Com a álgebra de feedbacks \javimos{21}{sensibilidade}:
$\Delta T_{2\times} \sim 2{,}5$--4\,K, e a maior incerteza mora nas
nuvens \javimos{6}{nuvens} — microfísica decidindo macroeconomia.

\section{A inércia: o oceano cobra depois}\label{sec:inercia}

\begin{formadif}[duas caixas (T2; Caso~2)]
\[ c_s\dot T_s = \Delta F - \frac{T_s}{\lambda} - \gamma(T_s - T_d),
\qquad c_d\dot T_d = \gamma(T_s - T_d). \]
Autovalores: $\sim$4 anos (camada de mistura — a escala do Pinatubo)
e $\sim$séculos (oceano profundo). Congelando o CO$_2$ hoje, o
planeta ainda aquece $\sim$0,4--0,6\,K: o \textbf{aquecimento
comprometido} — a dívida da inércia térmica \javimos{3}{capacidade
térmica da água}.
\end{formadif}

\section{O orçamento de carbono}\label{sec:tcre}

\quadro{Escrever a linearidade TCRE e deixar a turma desconfiada —
depois mostrar os dois efeitos não lineares que se cancelam (T4).}

\begin{noquadro}[TCRE (Caso~3)]
$\Delta T \simeq 0{,}45$\,K por 1000\,GtCO$_2$ acumuladas\\
emitidos até hoje: $\sim$2500\,GtCO$_2$ ($\to \sim$1,2\,K, confere
com o observado)\\
restam p/ 1,5\,°C: $\sim$800\,GtCO$_2$ $\approx$ 20 anos no ritmo
atual\\[2pt]
consequência: estabilizar $T$ $=$ emissão líquida \textbf{zero}
\end{noquadro}

A linearidade vem do log da forçante compensando a saturação dos
sumidouros (T4) — e transforma meta de temperatura em orçamento
finito, independente da trajetória (N3).

\section{Detecção, atribuição e assinaturas}\label{sec:detec}

\begin{formadif}[o sinal no ruído (Caso~4; N4)]
A variabilidade interna \javimos{21}{ENSO, vulcões} esconde a
tendência: com ruído realista, são necessárias $\sim$3--4 décadas
para 95\% de detecção — por isso o teste é sempre multidecadal, e uma
década fria isolada não refuta nada.
\end{formadif}

\begin{noquadro}[impressões digitais do mecanismo estufa (T5)]
troposfera aquece E estratosfera ESFRIA (Sol aqueceria as duas)\\
noites aquecem mais que dias; inverno mais que verão; Ártico mais que
tudo (gelo-albedo \javimos{21}{feedback})\\
OLR cai exatamente nas bandas do CO$_2$ (satélites, Cap.~8)
\end{noquadro}

\section{Impactos pelo curso}\label{sec:sintese}

Cada capítulo ganha seu pós-escrito: Clausius--Clapeyron intensifica
extremos de chuva ($+7\%$/K \javimos{4}{7\%/K}); ciclones tropicais
com teto MPI maior \javimos{16}{MPI $+2$\,K}; ondas de calor com
bulbo úmido no limite fisiológico \javimos{4}{$T_w$}; bloqueios e a
corrente de jato \javimos{11}{Rossby}; nível do mar (expansão térmica
$+$ gelo, T3); e a previsão sazonal-decadal como fronteira
\javimos{20}{ensembles}.

\section{Síntese do capítulo e do curso}

\begin{noquadro}[mapa final --- o fecho]
\[ \Delta F = 5{,}35\ln\frac{C}{C_0} \quad
\Delta T = \frac{\lambda_0\Delta F}{1-f} \quad
\tau_{oceano}: \text{anos e séculos} \quad
\Delta T \simeq \mathrm{TCRE}\times E_{cum} \]
o mesmo curso que prevê a chuva de amanhã fecha o balanço do século:
física, do barômetro ao planeta
\end{noquadro}

\section{Exercícios complementares}\label{sec:exercicios}

Soluções (professor) em \texttt{cap23\_solucoes.ipynb}.

\subsection*{Teóricos}

\begin{enumerate}[label=\textbf{T\arabic*.},itemsep=4pt]
  \item Explique a origem do $\ln(C/C_0)$ pela saturação do centro da
        banda e o avanço das asas; por que CFCs têm forçante linear?
  \item Obtenha as duas escalas de tempo do modelo de duas caixas e
        associe-as ao Pinatubo e ao aquecimento comprometido.
  \item Estime a subida do nível do mar por expansão térmica
        ($\alpha \simeq 2\times10^{-4}$\,K$^{-1}$) para 1\,K na
        camada de 700\,m, e compare com os $\sim$3,5\,mm/ano
        observados.
  \item Argumente a linearidade do TCRE pelos dois efeitos que se
        compensam e extraia a consequência ``líquido zero''.
  \item Explique por que a estratosfera esfria enquanto a superfície
        aquece, e por que isso descarta o Sol como causa.
\end{enumerate}

\subsection*{Numéricos (no notebook)}

\begin{enumerate}[label=\textbf{N\arabic*.},itemsep=4pt]
  \item Recupere o coeficiente $\sim$5,35 de um modelo de banda
        sintético e ache o regime onde o log falha.
  \item Mapeie o aquecimento comprometido no plano
        ($\lambda$, $\gamma$); que parâmetro domina?
  \item Rode três trajetórias com o mesmo acumulado e compare picos e
        caminhos de temperatura.
  \item Trace o ano de emergência do sinal contra a amplitude do
        ruído interno.
\end{enumerate}

\vspace{1em}
\noindent\rule{\linewidth}{0.4pt}\\
{\scriptsize Material derivado de R.~Stull, \emph{Practical Meteorology}
(CC BY-NC-SA 4.0); este documento herda a mesma licença.}

\end{document}
